% STATUS: rewritten 2026-11-05
% Sources read in full: McCrimmon 2004 (pp. 52, 57, 185, 360), Albert's
% exceptionality; Zel'manov 1979 read from the library copy 2026-11-09
% and cited first-hand; McCrimmon's pages are cited for his account, not as the source.
\section{Derivations of Jordan algebras}\label{sec:nonassociative:jordan-algebras}

Chapter~\ref{ch:nonassociative} opens with Jordan algebras, the historical
root of this chapter. Recall that a \emph{Jordan algebra}~\cite[Section~2.4, p.~57]{McCrimmon2004Taste} is a commutative
nonassociative algebra $J$ over a field $k$ of characteristic $\neq 2$
that satisfies the Jordan identity $(x\circ x)\circ(y\circ x) =
((x\circ x)\circ y)\circ x$. Albert showed that
$H_3(\mathbb{K})$ is a $27$-dimensional exceptional Jordan algebra
\cite[p.~971]{McCrimmon2004Taste}. Zel'manov later proved that no
other simple exceptional Jordan algebra exists, even in infinite
dimensions \cite[pp.~103--111]{Zelmanov1979Prime}.

A \emph{derivation} of a Jordan algebra $J$ has the same defining
property as in Chapter~\ref{ch:basic-theory}: a $k$-linear map
$\delta\colon J\to J$ with $\delta(xy)=\delta(x)y+x\delta(y)$. For an
associative algebra $A$ the inner derivations are the adjoint maps
$\operatorname{ad}(x)$, $x\in A$, and they exhaust the derivations in
nice cases. For a Jordan algebra a similar mechanism exists, but it
needs two indices: the commutators of two multiplication operators.

\begin{proposition}[{\cite[Problem~3.6~(3), p.~185]{McCrimmon2004Taste}}]
\label{prop:nonassociative:jordan-inner}
Let $J$ be a Jordan algebra over a ring in which $\tfrac{1}{2}$ exists,
and write $L_x$ for left multiplication by $x$. For each pair $x,y\in J$
the map
\[
D_{x,y}\;=\;4\,[L_x,L_y]\;=\;[V_x,V_y]
\]
is a derivation of $J$, where $[L_x,L_y]=L_xL_y-L_yL_x$ is the usual
commutator of multiplication operators.
\end{proposition}

Question~3.1 of McCrimmon~3.1 records that the derivations
$\mathrm{Der}(J)$ form a Lie algebra of endomorphisms under the bracket
$[D,E]=DE-ED$, with the inner derivations as an ideal
\cite[Question~3.1, p.~185]{McCrimmon2004Taste}. For special Jordan
algebras $J\subseteq A^{+}$ derivations come from the associative side:
any associative derivation of $A$ that leaves $J$ invariant is a Jordan
derivation of $J$ \cite[Problem~3.6~(2), p.~185]{McCrimmon2004Taste}.

The exceptional case gives the famous answer.

\begin{theorem}[{\cite[Historical remarks, p.~52]{McCrimmon2004Taste}}]
\label{thm:nonassociative:jordan-albert-F4}
Let $\mathcal{A}=H_3(\mathbb{K})$ be the Albert algebra over an
algebraically closed field of characteristic $0$, where $\mathbb{K}$
is the Cayley algebra. Then $\mathrm{Der}(\mathcal{A})$ is a
$52$-dimensional simple Lie algebra of type $F_4$, and
$\mathrm{Aut}(\mathcal{A})$ is the exceptional simple group of type
$F_4$.
\end{theorem}

McCrimmon records this as a discovery by Chevalley, Schafer, Freudenthal
and others in the early 1950s, with $G_2$ playing the analogous role for
the Cayley algebra itself \cite[p.~52]{McCrimmon2004Taste}. Zel'manov
proved that there is no simple exceptional Jordan algebra other
than the Albert algebras \cite[pp.~103--111]{Zelmanov1979Prime}. Together
with Theorem~\ref{thm:nonassociative:jordan-albert-F4} it singles out
$F_4$ as the only derivation algebra of an exceptional simple Jordan
algebra.

We close with the Jordan-variant built from a symmetric bilinear form.

\begin{example}[{\cite[Spin factor example, p.~360 and Problem~3.7~(1),
p.~185]{McCrimmon2004Taste}}]
\label{ex:nonassociative:jordan-spin}
Let $M$ be a $k$-module with a symmetric bilinear form $\sigma$. The
\emph{spin factor}
$JSpin(M,\sigma) = k\,\mathbf{1}\,\dot{+}\,M$ carries the Jordan
product
\[
(\alpha\,\mathbf{1}\,\dot{+}\,x)\bullet(\beta\,\mathbf{1}\,\dot{+}\,y)
\;=\;\bigl(\alpha\beta+\sigma(x,y)\bigr)\,\mathbf{1}\;\dot{+}\;
(\beta x+\alpha y).
\]
Its derivations $\delta$ satisfy $\delta(\mathbf{1})=0$ and
$\sigma(\delta(x),x)=0$ for all $x\in M$; thus $\delta|_{M}$ is
skew-symmetric for $\sigma$.
\end{example}

The Cayley algebra meets the Jordan side again in
Section~\ref{sec:nonassociative:alternative}: it is an alternative
algebra \cite[p.~148]{Baez2002Octonions}, and the next section studies
its derivations.

Neher classifies the Jordan triple systems whose derivation algebra or structure algebra is completely reducible, which is the derivation-algebra classification of the triple systems~\cite{WlJordanTripleSystemsWithComplet1984}.
